% 第 3 题三视图选择题题图组：2 个已知视图（正视图、俯视图，纵向叠放）+ 4 个选项侧（左）视图（横向排开）。
% 按固定源 PDF 矢量路径换算（1 pt → 0.03515 cm）；原点取图组左下（源页 (67.40, 566.30) pt），y 向上。
% 几何：俯视图为正方形 45.34×45.34 pt，其余为 1:2 竖长方形 45.34×90.67 pt（= 深 × 高）；两已知视图对角线均实线
% （正视图左上→右下、俯视图左下→右上）；选项 (A) 实线、(B) 虚线（均左上→右下），(C) 实线、(D) 虚线（均左下→右上）。
% 全组无尺寸标注；选项字母 (A)–(D) 在各选项图的左下方、图的左边缘之外（不在图下居中）。
\begin{tikzpicture}[line cap=round,line join=round,font=\small,x=0.03515cm,y=0.03515cm]
% ===== 已知视图 =====
% 正视图（竖长方形，实线对角线左上→右下）
\draw (133.30,186.28) rectangle (178.64,276.95);
\draw (133.30,276.95)--(178.64,186.28);
\node[anchor=base west] at (140.22,171.76) {正视图};
% 俯视图（正方形，实线对角线左下→右上）
\draw (133.30,118.27) rectangle (178.64,163.61);
\draw (133.30,118.27)--(178.64,163.61);
\node[anchor=base west] at (140.22,103.82) {俯视图};
% ===== 选项图（四个等大竖长方形，上下边齐平）=====
\draw (19.73,3.23) rectangle (65.07,93.91);
\draw (19.73,93.91)--(65.07,3.23);
\draw (104.30,3.23) rectangle (149.64,93.91);
\draw[dash pattern=on 2.99pt off 2.99pt] (104.30,93.91)--(149.64,3.23);
\draw (189.46,3.23) rectangle (234.79,93.91);
\draw (189.46,3.23)--(234.79,93.91);
\draw (274.90,3.23) rectangle (320.24,93.91);
\draw[dash pattern=on 2.99pt off 2.99pt] (274.90,3.23)--(320.24,93.91);
% 选项字母（图外左下）
\node[anchor=base west] at (0,3.03) {(A)};
\node[anchor=base west] at (85.01,3.03) {(B)};
\node[anchor=base west] at (170.02,3.03) {(C)};
\node[anchor=base west] at (255.02,3.03) {(D)};
\end{tikzpicture}
